% === Begin Label Data ===
% ID: 229
% 难度星级: 
% 标签: 
% 备注: 
% 组卷引用次数: 0
% === End  Label Data ===

\begin{problem}{2020}{G}{江苏卷}{20}{数列}
已知数列 $\{a_n\}(n\in\mathbb{N}^*)$ 的首项 $a_1=1$，前 $n$ 项和为 $S_n$．设 $\lambda$ 与 $k$ 是常数，若对一切正整数 $n$，均有 $\dfrac{1}{S_{n+1}}-\dfrac{1}{S_n^k}=\lambda a_{n+1}^k$ 成立，则称此数列为“$\lambda-k$”数列．

（1）若等差数列 $\{a_n\}$ 是“$\lambda-1$”数列，求 $\lambda$ 的值；

（2）若数列 $\{a_n\}$ 是“$\dfrac{\sqrt{3}}{3}-2$”数列，且 $a_n>0$，求 数列 $\{a_n\}$ 的通项公式；

（3）对于给定的 $\lambda$，是否存在三个不同的数列 $\{a_n\}$ 为“$\lambda-3$”数列，且 $a_n>0$？若存在，求 $\lambda$ 的取值范围；若不存在，说明理由．
\end{problem}